% ---------------------------------------------------------------------------
% OPTION A, CHAPTER 6 (RQ3). Written 2026-11-10 with AI assistance (Claude)
% from the drafted briefs in 05-implementation.tex (grounded tutor,
% anti-hallucination, prompt injection, assessment generation, review) and
% 06-evaluation.tex (citation gap), from docs/thesis/CHAPTER_MATERIAL.md §6
% and the verification note docs/thesis/verification-notes/2026-11-09-trust-rq3.md
% (branch claude/thesis-brief-verification-6573e8). Constants were re-checked
% against the code at the pin \thesisCommit on 2026-11-10. Labels of the moved
% sections are unchanged (sec:tutor, sec:retrieval, sec:anti-hallucination,
% sec:injection, sec:assessment, sec:review, sec:citation-gap,
% fig:tutor-sequence, fig:citation-gap); new labels carry a c6- prefix.
% The author must read, revise and own every sentence before submission.
% ---------------------------------------------------------------------------
% 2026-11-10 (cut to length): F12 and F16 moved to appendix B; the note on
% 24 June and the mechanisms table moved to appendix C.
\chapter{Trusting the Machine}
\label{ch:trust}

This chapter answers RQ3: how the quality of learning content written by a
language model can be verified and measured, and what happens when
verification fails. \Cref{ch:content} turned uploads into pages, and \cref{ch:loop} fed back to
the owner what readers do on them. Along the way a model writes three kinds
of text that a learner may take for true: an answer to a question, a quiz
question with its key, and the task behind a graded exercise.

One rule runs through the chapter. Where an answer can be computed, the
system computes it and does not ask the model. Where it cannot, a second model that did not write the content decides.
Where neither is possible, nothing checks the content, and the original page,
which a learner sees first, remains the reference. Most checks drop what
fails them; the few that only mark an item are named where they occur.

No check in the system decides whether a sentence of
generated prose is true. Every verifier below, that is, every step that computes, validates, filters,
or drops, tests something narrower: that a
page exists, that a phrase is on it, that a cited source was retrieved, that a
second model reaches the same key. \Cref{sec:c6-missing} lists what is
missing.

\section{The Page Stays the Reference}
\label{sec:c6-page}

The oldest trust decision in the project is not a verifier. From 24~February 2026 the original slide was shown beside the generated text
(commit \code{4c23d254}). The repository records no reason for it. In the reading view the uploaded page is the first of the three views of a
Lesson, and
the generated explanation and the plain text are one key away
(\code{LessonFacesScreen.tsx:152--163}). A learner who doubts a generated
sentence therefore always has the source in the same place. This costs
nothing at run time and it is the fallback for every mechanism that follows.

\section{Answers Grounded in the Space}
\label{sec:tutor}

A learner can ask in three places: the ask bar on the home screen, which
answers from a whole Space; a study session with the assistant, Luna; and
Explain and Ask on a single page of the reading view. \Cref{fig:groundedness} in
\cref{sec:rag} places them on a scale of how strongly the answer is tied to
the material. This section describes the strongest of them in detail and then
contrasts the other two.

\subsection{Retrieval Inside One Space}
\label{sec:retrieval}

The ask bar uses retrieval-augmented generation~\cite{lewis2020rag}. A vector
search (similarity above 0.65, at most 12 hits) and a full-text search (at
most 12 hits) run inside the Space, and reciprocal rank fusion with constant 60
keeps the best six (\code{ai/retrieval.py}). The scope is applied in the SQL
query, as a restriction to the Space's live Lessons, and not as a filter
afterwards (migration \code{20261014000000}).

This scope was wrong for a long time. Until 11~October 2026 the scoped search joined the embeddings on a column
that is empty for most rows, because embeddings are keyed by content hash;
the fix reached the main branch with PR~\#96 on 24~October, and when it
reached production is not recorded. Pages whose embeddings lacked that column
were invisible to the ask bar and to the global search; the fix, commit
\code{b66a97c6}, records 40 such rows out of 68 in one Space and adds a
database test that fails on the old join. How large the blind part was in other Spaces is not recorded; the
retrieval measurement of September, taken in another Space, found the right
page among its six hits for 26 of the 27 questions it could run there
(96\,\%) under the old join. Every later check depends on retrieval: a
refusal gate and a citation check are worth nothing when retrieval does not
see the material.

\subsection{Five Points of Control}
\label{sec:anti-hallucination}

\Cref{fig:tutor-sequence} in \cref{app:figures} shows the path of one
question.

\begin{enumerate}
  \item \emph{Scope.} The Space must be one the caller belongs to or owns
        (\code{v4\_spaces\_service.py}). Until PR~\#173 an owner who had not
        joined their own Space received a server error.
  \item \emph{Scoped retrieval,} as above.
  \item \emph{Input sanitising.} The learner's message is trimmed, cut at
        2,000 characters, and the delimiters with which the prompt fences uploaded text
        (\cref{sec:injection}) are rewritten so that the message cannot
        imitate them (\code{tutor.py:71--90}). The ask bar's own route refuses
        a question of more than 600 characters before that.
  \item \emph{Refusal gate.} If no retrieved page reaches similarity 0.65, a
        fixed ``not covered'' reply is returned before any model is called
        (\code{tutor.py:302--347}). The gate is a pure function of the
        retrieved list, so a test can assert on it without a model.
  \item \emph{Citation check.} The reply cites \code{[Source N]}. A citation
        whose number is not in the retrieved set is removed from the list of
        citations (\code{tutor.py:272--299}).
\end{enumerate}

The five points leave four gaps. None of them checks that the
content of the answer is supported by the cited page. None requires that an
answer cites anything. The invalid \texttt{[Source N]} tag is removed from the citation list
but stays in the text of the reply. The citation check compares against the
retrieved list and not against what fitted into the 10,000-character context,
so a source that was cut off still passes. In the terms of the attribution
literature~\cite{rashkin2023ais,gao2023alce}, the system verifies that a
citation points at retrieved material, not that the material supports the
claim.

\subsection{Strict First, Relaxed Later, Strict Again}
\label{sec:c6-history}

The first grounded tutor (commit \code{8d2eb095}, 2~May 2026) answered
strictly from the retrieved pages, cited them, and refused everything else.
On 24~June 2026 the refusal was removed from the tutor of the older interface (commit
\code{f39e93dc}). Since then that tutor may add knowledge from outside the
material if it puts it in a block quote and says that it goes
beyond the slides (\code{ai/prompts.py:506--540}). This is a prompt rule; no
code checks that the model keeps it.

Why the refusal was removed is not recorded;
\cref{sec:c6-history-detail} gives what the repository shows.

Two tutors therefore exist at \code{\thesisCommit}. The tutor of the older interface has no refusal gate and allows labelled
outside knowledge. The course tutor
behind the ask bar, which is the one learners reach in the current interface,
kept the gate and answers only from the Space. The system returned to strict answering when learners moved to the current
interface, not by a recorded decision.

\subsection{Luna and the Reading View: Checks on Structure}
\label{sec:c6-structure}

The other two surfaces do not use this retrieval and do not cite. Luna works
on the open page and the page before and after it, and can look up related
pages of the Space by full-text search. Explain and Ask get no search at all.
Both surfaces check the shape of what the model writes.

Luna names pages and lights up phrases on them. A move to a page that does not
exist is dropped, and a phrase survives only if it occurs literally on that
page (\code{luna/validate.py}; 12 unit tests). A check question must have
exactly one correct option. The file states the rule for Luna's structured output in its first line:
nothing the model writes reaches a learner unchecked, and what fails is
dropped, never shown broken.

Explain in the reading view sees the text of one page and the titles of up to four
neighbouring pages (\code{v4\_reader\_service.py:728}); Ask sees the same
page and up to four earlier turns of the conversation. Their output
must parse into a fixed set of blocks, and a block that does not is dropped
(\code{v4\_reader\_service.py}, \code{validate\_answer}). Beyond that they rely
on prompt rules: use only facts the page states, say when something is not
given, and label an invented example as hypothetical. These rules are
instructions to the model. Nothing verifies that the model follows them, and
this is the weakest of the three surfaces.

\subsection{Whether the Learner Sees the Source}
\label{sec:citation-gap}

A citation that is validated on the server and never shown gives the learner
nothing to check. That was the state at \code{\thesisEarlierCommit}
(\cref{fig:citation-gap} in \cref{app:figures}): the server computed and filtered citations, the one
component that rendered them was imported only by its own test, and the
screens in use either ignored the field or showed the page that happened to be
open. In the current interface the ask bar shows each citation as a link to
its page. The reading view's Explain and Ask still show none; their only source is
the open page itself. A trust mechanism that is not surfaced cannot be
perceived, which limits what any later study of perceived reliability could
find.

\section{Uploaded Text Is Not an Instruction}
\label{sec:injection}

Every prompt in the system carries text that the operator did not write. Any
signed-in user can create a Space, upload a PDF, and publish it, and any
signed-in user can join a published Space (\cref{sec:stakeholders}). A PDF can
therefore contain a sentence that is addressed to the model and not to the
reader. This is indirect prompt
injection~\cite{greshake2023indirect,owasp2025llm}. The audience of such a
document is the members of its Space.

The defence has two layers. The first is a rule in
every prompt that the page text and the learner's words are data and never
instructions. The second gives that rule a boundary that uploaded text cannot forge. Each prompt
fences uploaded text with markers, and a PDF that contains the same marker
could close the fence early. Since PR~\#179 one function rewrites exactly
those markers inside untrusted text into look-alikes that no fence matches
(\code{ai/untrusted.py}), in the tutors, in Luna, in the reading view, and in every
prompt of the parse. The idea is related to
spotlighting~\cite{hines2024spotlighting}. Before that change, at
\code{\thesisPreviousCommit}, retrieved text was inserted raw and the reading view's prompts lacked even
the rule of the first layer.

Rewriting the markers does not make injection impossible. The damage is bounded only where an output check exists:
Luna's literal phrase check, the citation check of the ask bar, and the block
validation of the reading view. An injected instruction that makes the reading
view's Explain write something false in well-formed blocks would pass. One earlier
hole widened the audience and is closed: the single-file upload did not check
who owned the course it was given (fixed in \code{722167b0}, with 10 tests).

\section{Questions That Carry a Key}
\label{sec:assessment}

A wrong sentence in an explanation can be checked against the page. A wrong
answer key is worse, because it marks a correct learner wrong and the learner
has no reason to doubt it. Most of the verification effort in the project went
here.

\subsection{From a Generated Question to a Stored One}
\label{sec:c6-lifecycle}

\thesisfigure{f24-question-lifecycle}
  {The life of one quiz question during a parse. Three stages can drop it, one only marks it, and a disputed key gets one
   repair on the same page; a second dispute drops the question.}
  {fig:question-lifecycle}

During a parse the model writes five to eight multiple-choice questions for
the whole deck from up to ten selected pages (\cref{sec:pipeline}). Each has
four options and names the one page that states its key. The prompt asks for
distractors that the page does not state as true and forbids negated
questions, two common rules of item
writing~\cite{haladyna2002mcq}. The code accepts at most twelve questions,
because the length of a model's list is not trusted either
(\code{parser/synthesis.py:52}). \Cref{fig:question-lifecycle} shows what
happens next.

\begin{description}
  \item[Resolve the key.] The model names the correct option as an index, a
        letter, or the option's text. If none of these matches an option, the
        question is dropped and logged. Until 16~June 2026 such a question was
        stored with option~A as its key (commit \code{f25dd41c}). The code
        states the reason: a wrong answer key is worse than fewer questions.
  \item[Structural gates.] No model is involved. A question is dropped if its
        cited page is missing, out of range, or empty, or if no key can be
        resolved (\code{parser/quiz\_grounding.py:127}). A number in the
        key that is not on the page, or a distractor that the page supports
        better than the key, only marks the question, because the switch that would
        make these checks drop it is off (\code{parser/synthesis.py:55--63}).
        A marked question is still shown to learners; the mark is stored with
        it and counted in the parse receipt.
  \item[Lexical support.] A score measures how much of the key's wording
        occurs on the cited page; below 0.5 the question is marked. It is
        not dropped, for the reason measured in \cref{sec:c6-evaluation}.
  \item[Answer check.] A second model reads the cited page, the question, and
        the options, without the key, and answers as a careful student would
        (PR~\#166). If it picks another option, finds a second defensible
        option, or finds none, the key is disputed.
  \item[Repair or drop.] A disputed question is rewritten once against the
        same page and checked again. A second dispute drops it. A repair may
        not move the question to another page.
\end{description}

Two properties of the answer check are design decisions. First, the model that checks a question is not the model that wrote it. The question is written by the small model at the head of the
bulk chain; the check runs on the quality chain, led by a larger reasoning
model (\cref{sec:provider-strategy}). At the pin the heads of both chains are models of the
same provider (\code{orchestrator.py:389}, \code{:393}), so the
independence is one of model, not of vendor. The reason for using a different model is recorded with a measurement: when a
model of the generating model's own family checked the question that started
this work, an ambiguous question whose options all appear on its page, it
agreed with the wrong key twice. That a
model tends to favour text like its own is also reported in the
literature~\cite{panickssery2024selfpreference}. Second, an outage and a dispute are treated differently. If the check itself fails or its reply
cannot be read, the question is kept, so that a provider outage never empties
a quiz. If the check disputes the key, the question is rewritten once and dropped only
if the rewrite is disputed too. The outcome of
every stage is counted in the parse receipt of the Lesson.

\subsection{Grading by Computation}
\label{sec:c6-grading}

The rehearsal, a short mixed test that a model writes from a Space's pages,
and the reading view's checkpoints, which a model writes for the pages of one
idea and which are not the deck quiz of \cref{sec:c6-lifecycle}, use five
kinds of question.
Three of them are about functional dependencies in the database course the
corpus comes from: which dependencies hold on a given table, the closure of a
set of attributes, and which dependencies violate Boyce--Codd normal form (BCNF). For these the model
writes the task but not the verdict. The answer is computed from the
question's own data, on the client and again on the server
(\code{rehearsal/grade.ts}, \code{v4\_rehearsal\_service.py:43--105}), and a
task whose computed answer is trivial or empty is refused before a learner
sees it (PR~\#75). The file says it in one sentence: the model does not get
the last word on the answers.

The other two kinds cannot be computed. A multiple-choice question is checked
for shape (exactly one option marked correct) and an ordering question for
distinct items, and both are then graded against the model's key. Computation
covers exactly the questions that have a formal answer. It is the strongest
verifier in the system and the narrowest.

\section{Measuring the Verifiers}
\label{sec:c6-evaluation}

A verifier can fail in two directions. It can drop a good question, which
costs content, or keep a bad one, which costs trust. We measured both for
every verifier of answer keys on the same fixed probe.

\subsection{Method}
\label{sec:c6-method}

The probe is 16 pairs, written by hand on 11~October 2026 before any verifier
was scored on them (\code{backend/eval/data/grounding\_probe.py}). Eight pairs
are in English, from a lecture on transactions, and eight in German, from a
lecture on storage management. Both decks are disjoint from the decks on
which any threshold was fitted. Each pair holds one correct key that
paraphrases the page and one wrong key that stays on the topic of the page.
The correct keys were written under a rule that forbids copying three words
in a row from the page, and the wrong keys are plainly false
(\code{docs/DECK\_LAB\_2026\_10.md}). A
verifier that works on wording will fail the first; a verifier that works on
topic will fail the second. The expected shape was that the lexical score
drops paraphrases, and that a semantic score repairs this without keeping the
wrong keys.

For the semantic score the procedure was fixed in writing before any embedding
was requested: the scoring rule; that the threshold is fitted only on a
separate calibration set of key pairs from other decks; that the probe is
scored once; the acceptance criterion (at most two
calibration errors, at most two good keys dropped, at most one wrong key kept);
what follows from a pass and from a fail; and a named risk, that the question
text without its options (the stem) could dominate the embedding and raise
wrong keys above the threshold. The threshold itself,
0.774, was not registered in advance. It is the fitted value.

\subsection{Results}
\label{sec:c6-results}

\begin{table}[htbp]
  \centering
  \small
  \caption{Three verifiers of answer keys on the same 16 pairs: correct keys
    dropped and wrong keys kept. The acceptance criterion for enforcement was at most two good keys dropped
    and at most one wrong key kept. For the answer check, counts are over the
    calls that returned a verdict; one call did not. A count of 0 of 16 is compatible with an error rate of up to
    21\,\% (exact 95\,\% interval). Sources: \code{parser/synthesis.py:55--80}; PR~\#94 and tag
    \code{archive/claude/arc8-semantic-grounding}; PR~\#166.}
  \label{tab:verifiers}
  \begin{tabular}{@{}p{0.33\textwidth}ccp{0.22\textwidth}@{}}
    \toprule
    Verifier & Good dropped & Wrong kept & Consequence \\
    \midrule
    Lexical score, threshold 0.5 & 14 of 16 & 0 of 16 & marks only \\
    Similarity score, fitted threshold 0.774 & 4 of 16 & 8 of 16 &
      criterion failed; not merged \\
    Answer check & 0 of 16 & 0 of 16 & enforced \\
    \bottomrule
  \end{tabular}
\end{table}

\Cref{tab:verifiers} gives the result. The lexical score never keeps a wrong
key, but it drops 14 of 16 correct paraphrases, among them all eight German
ones. Enforced, it would delete most good questions in a German Lesson. It
was therefore set to mark and not to drop, and the constant carries a comment
that forbids enforcing it without a new measurement.

The similarity score failed its acceptance criterion, and it failed in the
direction named in advance as its risk.
It dropped 4 good keys and kept 8 of 16 wrong ones, six of them German,
against a limit of one. On the calibration set it made ten errors against a
limit of two. On the probe, similarities fell into a narrow band between
0.70 and 0.91,
because the stem dominates the embedding and a wrong key on the same topic
scores like a right one. The consequence had been fixed in advance: the
wiring into the parse was reverted, and tuning the threshold against the same
probe was ruled out. The experiment was closed without merging (PR~\#94); its
code and record exist only under an archive tag. This negative result is
the reason the next attempt did not use similarity at all.

The answer check dropped none of the 16 good keys and kept none of the 16
wrong ones. On the question that started the work, which a careful student
could answer differently from its key, it disputed the key in five of five
runs and accepted a corrected version of the same question in five of five. The run was made on
30~October 2026 with the reasoning model at the head of the quality chain
(PR~\#166). Unlike the similarity score, the answer check was not held to an acceptance
criterion fixed in advance. It was built for that one question and measured in
the same change, on the probe that already existed (commit
\code{461afc40}), so its result is in-sample. Before the check itself is called validated, it should be scored once on new
pairs from other decks against an acceptance criterion fixed first.

One call of the run returned no verdict, because the provider limited the
rate and the fallback timed out. The script counts such a call as an error
(\code{backend/\allowbreak eval/\allowbreak quiz\_answer\_check.py:101--103}),
and the record of the run does not say which pair it was. The deployed check
would have kept that question whichever key it carried, so
\cref{tab:verifiers} describes the check when it answers.

The refusal gate was measured separately at \code{\thesisCommit},
in November 2026, with expectations committed before the run
(\code{docs/thesis/experiments/retrieval-rerun/}). In one Space of three
Lessons and 201~pages we asked 27 questions that have an answer
there, drafted with AI assistance from the page text and reviewed by the
author, and 30 that have none: ten far from the subject, ten from computer
science outside databases, and ten on database topics that the three Lessons
do not cover. With the settings of production the gate let all 27 through and
refused all 30. The 27 answerable questions had already passed the gate in
September, so this half replicates a known outcome; the new evidence is the
30 unanswerable ones. Zero errors in 30 is compatible with a rate of wrongly
accepted questions of up to 12\,\%. The similarities of the answerable and the unanswerable questions do not
overlap, but the gap is narrow: no unanswerable question reaches more than
0.634, the weakest answerable question reaches 0.659, and the threshold of
0.65 lies between them. A
threshold of 0.75 would have refused 13 of the 27 answerable questions. We had
predicted that some of the absent database topics would pass the gate; none
did. For the answerable questions, the expected page was among the first five hits in 25 of 27 cases (the
closest setting the grid measured to the six hits the ask bar keeps).

\subsection{What the Numbers Do Not Show}
\label{sec:c6-threats}

The probe is small and was written by the developer of the system, in two
subjects from one university. Sixteen pairs show how a verifier behaves on
these chosen cases; they do not estimate an error rate in production.
The wrong keys are on the topic but plainly false, so a keep rate of zero
shows that such keys are rejected, not that subtler errors are. Because the
correct keys avoid the page's wording by rule, the lexical result is partly
a product of the probe. The failure that started the answer check, an
ambiguous question whose options all appear on the page, is represented by
one case.
The answer check was measured once, on one day, with one model at the head of
its chain, and a model behind an API can change without notice. No measurement
covers the prose of Explain, Ask, or Luna. The refusal gate was measured in
one Space, with 57 English questions written
by the developer with AI assistance, ten per group of unanswerable questions.
Their wording may suit retrieval better than a student's would, and a margin
of 0.025 on this sample does not show how the threshold behaves on other
material.
A question that passes the gate is also not thereby answered correctly: what
the model then writes was not measured.

The answer check is itself a model call, so what happens when verification
fails also depends on what happens when a provider fails. Property tests
written for that case, and two cost defects that planning them uncovered,
are reported in \cref{sec:c6-cost}.

\section{Where Nobody Checks}
\label{sec:c6-missing}

\subsection{No Human Approval}
\label{sec:review}

No person has to approve generated content before learners see it. In the older interface a professor can review a batch of uploads, but a lecture has no
draft state: its content is visible as soon as the parse ends, and ``done
reviewing'' is stored only in the professor's browser
(\code{BatchReviewPage.tsx}). A slide that was corrected by hand keeps its review mark. In the current interface the owner of a Space adds Lessons with no
review screen at all, and the members of a published Space see a new Lesson
the moment its parse creates it. Publishing controls who can find and join a
Space, not what its members see.

This is a consequence of a decision made for another reason. The reading view was built so that a page can be read seconds after the
upload (\cref{sec:pipeline}).
A real approval step would need a draft state on the Lesson and would delay exactly
that. The automatic checks of this chapter are therefore the only line
between the model and the learner, which is why an item that fails a check is dropped rather than passed on for
a human to catch.

\subsection{The Mechanisms in One Table}
\label{sec:c6-inventory}

\Cref{tab:mechanisms} in \cref{sec:c6-inventory-table} groups the
26~mechanisms that exist in the code by what they guard. Most verifiers guard
structure and answer keys, where a check can be exact. Meaning is guarded by prompt rules on every surface and, for quiz keys only,
by the answer check.

Three things do not exist. There is no entailment check: nothing tests directly whether a page logically
supports a key or a sentence. The answer check tests a key only indirectly,
by whether a second model reaches it, and the lexical module says itself
that it is no such test. There is no detector for invented content in the prose of
Explain, Ask, or Luna; current surveys treat this as an open
problem~\cite{ji2023hallucination,huang2025hallucination}. And nothing checks
that the older tutor really puts outside knowledge into a block quote.

\section{Answer to RQ3}
\label{sec:c6-answer}

The measurements in this chapter cover answer keys and the refusal gate.
The prose of explanations and answers is neither verified nor measured.

\emph{How can the quality of content written by a language model be verified?}
In this system, in three ways of falling strength. By computation, where the
content has a formal answer: three kinds of task are graded by code and the
model's key is never used. By the answer check, for the quiz questions written during parsing: a different model answers each question
from its cited page without the key, and if that check is unavailable the
question is kept. By structure everywhere else: pages must exist, phrases
must be on the page, citations must point at retrieved sources, blocks must
parse. This weakest kind also covers the multiple-choice and ordering
questions of the reading view and the rehearsal, which are graded against the
model's key.

\emph{How can it be measured?} With a small probe that is written before the
verifier is scored and that contains both kinds of error. On 16 pairs the
lexical score dropped 14 good keys, the similarity score kept 8 wrong ones
and failed an acceptance criterion that had been fixed in advance, and the
answer check made neither error in the calls that returned a verdict. The refusal gate, tested in one Space with 57 questions,
let through all 27 that have an answer and refused all 30 that have none, by
a narrow margin. The failure of the similarity score shaped the design as much as the success
of the answer check.

\emph{What happens when verification fails?} The item is dropped: a key that
cannot be resolved, a question whose key is disputed twice, a move to a page
that does not exist, a block that does not parse. A question outside the
material is refused before any model is called. When the verifier itself is
unavailable the system keeps the content, by decision, so that an outage does
not remove it. And where no verifier exists, for the meaning of generated
prose, the original page remains the first thing a learner sees.
